\chapter{Results}
\label{ch:results}

\begin{itemize}
    \item Representative scenario: AoS baseline with sorting
    \begin{itemize}
        \item Runtime by variant.
        \item Speedup by variant.
        \item Runtime decomposition by variant.
        \item L2 miss ratio / miss rate by variant.
        \item Stall rates by variant.
        \item L3 data volume / bandwidth by variant.
        \item Identify whether sorting improves locality.
        \item Identify whether locality improvement translates into runtime improvement.
    \end{itemize}

    \item Sorting resolution comparison
    \begin{itemize}
        \item Cell resolution
        \begin{itemize}
            \item Expected low overhead.
            \item Coarse locality improvement.
            \item Likely strongest practical baseline.
        \end{itemize}
        \item Block resolution
        \begin{itemize}
            \item Expected lower overhead than particle resolution.
            \item May preserve larger spatial chunks.
            \item Interpret together with block size.
        \end{itemize}
        \item Particle resolution
        \begin{itemize}
            \item Expected best spatial order.
            \item Highest overhead.
            \item Runtime may be dominated by key generation and sorting.
        \end{itemize}
    \end{itemize}

    \item Sorting order comparison
    \begin{itemize}
        \item Linear
        \begin{itemize}
            \item Lowest overhead.
            \item If it wins, overhead dominates over sophisticated locality.
        \end{itemize}
        \item Morton
        \begin{itemize}
            \item Middle ground between locality and overhead.
        \end{itemize}
        \item Hilbert
        \begin{itemize}
            \item Strong locality expectation.
            \item Must justify overhead through measured gains.
        \end{itemize}
    \end{itemize}

    \item SoA-based traversal results
    \begin{itemize}
        \item Compare AoS and SoA configurations.
        \item Determine whether sorted SoA order improves force traversal.
        \item Check whether SoA overhead is offset by traversal improvements.
        \item Interpret runtime decomposition carefully.
    \end{itemize}

    \item Neighbor-list sorting results
    \begin{itemize}
        \item Compare with and without neighbor-list sorting.
        \item Determine whether neighbor access order improves.
        \item Check whether rebuild overhead outweighs traversal benefit.
    \end{itemize}

    \item Scaling with particle count
    \begin{itemize}
        \item Fixed-density increasing domain size.
        \item Runtime scaling.
        \item Speedup scaling.
        \item Locality metric scaling.
        \item Determine whether sorting becomes more useful for larger working sets.
    \end{itemize}

    \item Scaling with density
    \begin{itemize}
        \item Fixed-domain increasing density.
        \item Longer neighbor lists.
        \item More force traversal pressure.
        \item Determine whether locality effects become stronger with denser configurations.
    \end{itemize}
\end{itemize}
